\documentclass[11pt]{article}

\usepackage[margin=1in]{geometry}
\usepackage[T1]{fontenc}
\usepackage[utf8]{inputenc}
\usepackage{setspace}
\usepackage{amsmath}
\usepackage{booktabs}
\usepackage{array}
\usepackage{xcolor}
\usepackage{hyperref}
\usepackage{listings}
\usepackage{longtable}
\usepackage{enumitem}
\usepackage{caption}

\onehalfspacing
\hypersetup{colorlinks=true, linkcolor=blue, urlcolor=blue}
\setlist[itemize]{topsep=4pt,itemsep=4pt}
\setlist[enumerate]{topsep=4pt,itemsep=4pt}

\definecolor{codegray}{RGB}{248,248,248}
\definecolor{darkblue}{RGB}{20,60,120}

\lstdefinestyle{pythonstyle}{
  language=Python,
  basicstyle=\ttfamily\small,
  backgroundcolor=\color{codegray},
  keywordstyle=\color{darkblue}\bfseries,
  commentstyle=\color{gray}\itshape,
  stringstyle=\color{teal},
  showstringspaces=false,
  breaklines=true,
  frame=single,
  columns=fullflexible
}

\lstdefinestyle{textstyle}{
  basicstyle=\ttfamily\small,
  backgroundcolor=\color{codegray},
  showstringspaces=false,
  breaklines=true,
  frame=single,
  columns=fullflexible
}

\title{SOEN 6011 Delivery 2 Report\\F3: Hyperbolic Sine Function $\sinh(x)$}
\author{Arvind Lakshmanan\\Student ID: 40310757\\Concordia University}
\date{Summer 2026}

\begin{document}
\maketitle

\section*{Overview}

This document gives the Delivery 2 answer for the assigned function F3, $\sinh(x)$. The D1 implementation used the direct exponential definition with a textual interface. For D2, the implementation is modified to satisfy the scratch implementation requirement, the Tkinter GUI requirement, exception handling, repository documentation, and updated requirements.

\section{Problem 5: From-scratch Python implementation with Tkinter GUI}

\subsection{D2 implementation decision}

The selected D2 algorithm is the Maclaurin series for the hyperbolic sine function:

\[
\sinh(x)=x+\frac{x^3}{3!}+\frac{x^5}{5!}+\frac{x^7}{7!}+\cdots
\]

This algorithm was selected for D2 because it can be implemented from scratch using only arithmetic operations. The implementation does not use \texttt{math.sinh}, \texttt{math.exp}, \texttt{math.factorial}, NumPy, SciPy, or other mathematical library functions. Tkinter is used for the graphical user interface only.

\subsection{Why the D1 implementation was modified}

The D1 implementation used the formula

\[
\sinh(x)=\frac{e^x-e^{-x}}{2}.
\]

That was suitable for D1 because D1 required a textual interface and did not yet require a scratch implementation. In D2, built-in or library mathematical functions are not used for calculating the function. Therefore, the D2 implementation replaces the direct exponential method with a from-scratch Maclaurin method.

\subsection{Subordinate functions implemented}

The implementation includes subordinate functions that support the main function calculation:

\begin{itemize}
    \item \texttt{absolute\_value(value)}: returns a non-negative value without using \texttt{abs()}.
    \item \texttt{is\_nan(value)}: detects NaN using the property that NaN is not equal to itself.
    \item \texttt{is\_infinite\_or\_too\_large(value)}: rejects infinity and unsafe huge values.
    \item \texttt{parse\_supported\_real\_number(text)}: converts GUI input to a finite supported real number.
    \item \texttt{calculate\_sinh\_from\_scratch(x\_value)}: computes $\sinh(x)$ using a recurrence form of the Maclaurin series.
\end{itemize}

\subsection{Series recurrence}

Instead of recomputing powers and factorials from the beginning, the implementation uses the following recurrence:

\[
\text{term}_0=x
\]

\[
\text{term}_n=\text{term}_{n-1}\times \frac{x^2}{(2n)(2n+1)}
\]

The running sum is updated by adding each term. The loop stops when the absolute value of the current term is smaller than the tolerance condition:

\[
|\text{term}| \leq \text{TOLERANCE}\times(1+|\text{sum}|).
\]

\subsection{Input range}

The supported D2 range is:

\[
-20 \leq x \leq 20.
\]

This range is used because the scratch series implementation is accurate and responsive for ordinary calculator inputs in this interval. Larger ranges can be added later using more advanced numerical techniques, such as argument reduction.

\subsection{Exception handling and helpful messages}

The source code defines custom exception classes:

\begin{itemize}
    \item \texttt{SinhInputError}: used for blank input, non-numeric input, NaN, infinity, and out-of-range values.
    \item \texttt{SinhConvergenceError}: used if the series does not converge within the maximum iteration limit.
\end{itemize}

The GUI catches these exceptions and prints messages that explain what went wrong and how the user can correct the input.

\subsection{Graphical user interface}

The GUI is implemented using Tkinter. It includes:

\begin{itemize}
    \item an input box for $x$,
    \item a Calculate button,
    \item a Clear button,
    \item an Exit button,
    \item a result area showing the input, the calculated value, algorithm name, and number of terms used,
    \item a status message for success or error conditions.
\end{itemize}

\subsection{Manual validation examples}

\begin{center}
\begin{tabular}{r r l}
\toprule
\textbf{Input} & \textbf{Approximate output} & \textbf{Purpose} \\
\midrule
0 & 0 & zero case \\
1 & 1.1752011936438 & common positive input \\
-1 & -1.1752011936438 & common negative input \\
3 & 10.0178749274099 & larger positive input \\
-3 & -10.0178749274099 & larger negative input \\
20 & 242582597.704895 & upper boundary \\
\bottomrule
\end{tabular}
\end{center}

\subsection{Problem 5 CASTROFF prompt record}

\textbf{Prompt type:} implementation-generation prompt with constraints and verification criteria.

\begin{lstlisting}[style=textstyle]
Context: I am preparing SOEN 6011 Delivery 2 for F3, sinh(x). D1 used a textual Python implementation. D2 requires modifying it to implement the function from scratch in Python and adding a Tkinter GUI.
Audience: A course evaluator who will check that no Python math-library function is used for the function calculation.
Specific task: Generate a Python design and source code structure for sinh(x) using the Maclaurin series, with Tkinter input/output, custom exceptions, helpful error messages, and no math.sinh, math.exp, math.factorial, NumPy, or SciPy.
Tone: Clear, academic, and implementation-focused.
Requirements: Support finite real input, reject invalid input, explain the range, and separate parsing, calculation, and GUI responsibilities.
Output format: Provide code and a short explanation of the algorithm and exception handling.
Feedback loop: After producing the output, check whether each D2 constraint is satisfied.
Final check: Identify any built-in or library mathematical function that should be removed.
\end{lstlisting}

\textbf{Output use:} The output was used to structure the code, separate GUI and calculation responsibilities, and document why the Maclaurin series is appropriate for a scratch implementation. The final source code was manually reviewed and modified.

\section{Problem 6: Public distributed version control system}

The source code must be placed in a public distributed version control system. The intended hosting service is GitHub.

\subsection{Repository address}

\textbf{Repository URL to include after upload:}

\begin{lstlisting}[style=textstyle]
https://github.com/<your-username>/soen6011-d2-f3-sinh
\end{lstlisting}

Before final submission, replace the placeholder URL with the actual public repository address.

\subsection{Repository contents}

The repository should contain:

\begin{itemize}
    \item \texttt{D2\_F3\_sinh\_gui.py}
    \item \texttt{README.md}
    \item \texttt{Arvind\_Lakshmanan\_40310757\_D2\_Report.pdf}
    \item \texttt{Arvind\_Lakshmanan\_40310757\_D2\_Slides.pdf}
    \item \texttt{commit\_messages.md}
\end{itemize}

\subsection{High-quality commit messages}

High-quality commit messages should be specific, action-oriented, and connected to meaningful project changes. Suggested commit sequence:

\begin{enumerate}
    \item \texttt{docs: add D2 README with setup and run instructions}
    \item \texttt{feat: implement from-scratch Maclaurin series for sinh}
    \item \texttt{feat: add Tkinter GUI with input validation and helpful errors}
    \item \texttt{refactor: separate parsing, calculation, and GUI responsibilities}
    \item \texttt{docs: update D2 requirements and traceability}
\end{enumerate}

\subsection{README content}

The README file explains the project purpose, supported input range, run command, file list, example outputs, and the requirement to replace the placeholder GitHub URL with the actual public repository URL.

\subsection{Problem 6 CASTROFF prompt record}

\textbf{Prompt type:} documentation and version-control planning prompt.

\begin{lstlisting}[style=textstyle]
Context: I need to prepare the D2 repository evidence for a SOEN 6011 project. The project is F3, sinh(x), implemented from scratch with a Tkinter GUI.
Audience: A professor or TA checking repository quality.
Specific task: Draft README content and high-quality Git commit messages. The README must explain how to run the program, supported input, files, and the from-scratch constraint.
Tone: Professional and concise.
Requirements: Include a placeholder for the public GitHub address and a reminder to replace it before submission.
Output format: Markdown README and a list of commit messages.
Feedback loop: Check whether the README includes purpose, run command, input range, and file list.
Final check: Ensure no claim says the repository already exists unless the public URL has been added.
\end{lstlisting}

\textbf{Output use:} The output was used to prepare \texttt{README.md} and \texttt{commit\_messages.md}. The repository URL remains a placeholder until the project is uploaded to a public repository.

\section{Problem 7: Updated D2 requirements}

The D1 requirements are modified to account for the D2 scratch implementation and Tkinter GUI.

\subsection{Definitions for D2}

\begin{longtable}{p{0.27\textwidth}p{0.65\textwidth}}
\toprule
\textbf{Term} & \textbf{D2 definition} \\
\midrule
\endhead
From scratch & The function calculation must not use Python built-in or library mathematical functions such as \texttt{math.sinh}, \texttt{math.exp}, or \texttt{math.factorial}. \\
GUI & A graphical user interface created using Tkinter for entering input and viewing output. \\
Supported real input & One finite real number that can be represented as a Python floating-point value and lies in the interval $[-20,20]$. \\
Helpful error message & A message that identifies the problem and guides the user toward a valid input. \\
Convergence & The Maclaurin series reaches the stopping condition before the maximum number of terms. \\
\bottomrule
\end{longtable}

\subsection{Updated requirements list}

\begin{longtable}{p{0.18\textwidth}p{0.55\textwidth}p{0.14\textwidth}p{0.08\textwidth}}
\toprule
\textbf{ID} & \textbf{Updated requirement} & \textbf{Rationale} & \textbf{Priority} \\
\midrule
\endhead
REQ-F3-D2-01 & The system shall provide a Tkinter graphical user interface for input and output. & D2 requires GUI support. & High \\
REQ-F3-D2-02 & The system shall accept exactly one finite real number $x$ through the GUI. & Prevents ambiguous input. & High \\
REQ-F3-D2-03 & The system shall support inputs in the range $-20 \leq x \leq 20$. & Maintains reliable performance for the scratch series. & High \\
REQ-F3-D2-04 & The system shall calculate $\sinh(x)$ using a from-scratch Maclaurin series implementation. & Satisfies the scratch requirement. & High \\
REQ-F3-D2-05 & The system shall not use \texttt{math.sinh}, \texttt{math.exp}, \texttt{math.factorial}, NumPy, SciPy, or equivalent mathematical library functions for the calculation. & Prevents prohibited shortcuts. & High \\
REQ-F3-D2-06 & The system shall stop the series when the term is smaller than the defined tolerance condition. & Supports controlled approximation. & High \\
REQ-F3-D2-07 & The system shall raise and handle a helpful input error for blank, non-numeric, NaN, infinity, and out-of-range inputs. & Improves user experience. & High \\
REQ-F3-D2-08 & The system shall raise and handle a convergence error if the required tolerance is not reached within the maximum term count. & Handles computational failure safely. & Medium \\
REQ-F3-D2-09 & The system shall display the input, calculated output, algorithm name, and number of terms used. & Gives transparent feedback to the user. & Medium \\
REQ-F3-D2-10 & The system shall separate parsing, calculation, and GUI responsibilities in the source code. & Improves maintainability. & Medium \\
REQ-F3-D2-11 & The source code shall run using standard Python without depending on a specific IDE. & Satisfies portability expectation. & High \\
REQ-F3-D2-12 & The source code shall be placed in a public distributed version control system with a README file. & Satisfies D2 repository requirement. & High \\
\bottomrule
\end{longtable}

\subsection{Traceability from D1 to D2}

\begin{center}
\begin{tabular}{p{0.22\textwidth}p{0.34\textwidth}p{0.34\textwidth}}
\toprule
\textbf{D1 concern} & \textbf{D2 modification} & \textbf{Reason} \\
\midrule
Textual interface & Replaced by Tkinter GUI & D2 requires GUI. \\
Direct exponential formula & Replaced by Maclaurin series & D2 requires scratch implementation. \\
Use of \texttt{math.exp} & Removed from function calculation & Library math functions are not allowed for scratch calculation. \\
Input error handling & Expanded with custom exceptions & D2 requires exception handling and helpful messages. \\
Supported range $[-709,709]$ & Updated to $[-20,20]$ & Ensures reliable scratch-series performance. \\
\bottomrule
\end{tabular}
\end{center}

\subsection{Problem 7 CASTROFF prompt record}

\textbf{Prompt type:} requirements-revision prompt.

\begin{lstlisting}[style=textstyle]
Context: D1 requirements were written for a textual sinh(x) calculator using a direct exponential formula. D2 changes the implementation to a from-scratch Maclaurin series and a Tkinter GUI.
Audience: SOEN 6011 evaluator checking requirements traceability.
Specific task: Update the D1 requirements for D2. Requirements must cover GUI input/output, scratch calculation, input range, helpful exceptions, no math libraries, README, and public version control.
Tone: Formal requirements wording using shall statements.
Requirements: Provide unique identifiers, rationale, and priorities. Make implementation decisions explicit.
Output format: A requirements table and a traceability table from D1 to D2.
Feedback loop: Check whether every D2 Problem 5 and Problem 6 constraint is represented by at least one requirement.
Final check: Avoid unsupported claims and keep assumptions explicit.
\end{lstlisting}

\textbf{Output use:} The output was used to produce the updated requirements table and the D1-to-D2 traceability table. The wording was manually reviewed to keep each requirement verifiable.

\section{Appendix: Python source code}

\lstinputlisting[style=pythonstyle]{D2_F3_sinh_gui.py}

\end{document}
